\documentclass[10pt,a4paper]{amsart}
\usepackage[margin=19mm]{geometry}
\usepackage{amsmath,amssymb,mathtools}
\usepackage[scr=rsfs,cal=euler]{mathalfa}
\usepackage{xfrac}
\usepackage[unicode,hidelinks,bookmarks=false]{hyperref}
\newcommand{\ie}{i.e.\ }
\newcommand{\cf}{cf.\ }
\newcommand{\resp}{respectively\ }
\newcommand{\Subsection}{\subsection}
\providecommand{\nobreakdash}{\nobreak\hbox{-}}
\setlength{\emergencystretch}{2em}
\multlinegap0pt
\usepackage{amssymb}
\usepackage{mathtools}
\usepackage{enumerate}
\newtheorem{lem}{Lemma}
\DeclareMathOperator{\td}{d\!}
\DeclareMathOperator{\te}{e}
\title{The decrease in sectional curvature of the Fisher--Rao manifold of beta distributions}
\author{Feng Qi}
\date{Source: arXiv:2609.10628v1 (CC BY 4.0)}
\begin{document}
\maketitle

\noindent\emph{Source and licence.} The decrease in sectional curvature of the Fisher--Rao manifold of beta distributions. Version arXiv:2609.10628v1 (CC BY 4.0), distributed by arXiv under the Creative Commons Attribution 4.0 International licence, http://creativecommons.org/licenses/by/4.0/, as recorded in the arXiv licence metadata for this exact version and shown on the abstract page https://arxiv.org/abs/2609.10628v1. Full licence notice: \texttt{licenses/arxiv-2609.10628-CC-BY-4.0.txt}.\par\smallskip

\noindent\textit{Editorial scope.} Counted unit: the article's lem lem:general (source0.tex lines 111--121). The article's complete original proof of it is delivered with it (source0.tex lines 123--238, one \texttt{proof} environment of 116 lines). No mathematics is written, repaired or re-derived: the statement and the proof are the article's own lines, in the article's own notation, and no retained line is substituted or re-flowed.  No further source-local context is needed: every cross-reference of every retained line resolves inside the two delivered spans.  Nothing was omitted from any retained span, so the package carries no omission record.  No retained line contains a citation, so the wrapper carries no bibliography and no bibliography entry.  No retained line contains a figure environment, an image-inclusion command or drawing code, so no image is delivered and the package declares no asset dependency.  Editorial formatting only, in addition: an AMS \texttt{amsart} preamble that restates the article's own declarations and macros used by the retained lines, namely the environments \texttt{lem} on separate counters, together with the standard packages those lines need.  The retained environments print in this document as Lemma 1 in order of appearance, whereas the article prints the same items with the numbering of its own full document.  The complete native source of the article as distributed in the arXiv e-print for this exact version is bundled unchanged as \texttt{sources/209937/source0.tex}.
\begin{lem}\label{lem:general}
Let $f(x)$ be thrice continuously differentiable and $H(x)=-\frac{f(x)}{f'(x)}$ on $(0,\infty)$. If
$$
f(x)>0,\quad f'(x)<0, \quad 0<H'(x)<1, \quad H''(x)\geq 0
$$
on $(0,\infty)$, then, for every given $a>0$, the function
$$
F_a(x)=\frac{1}{1-\frac{f(x+a)}{f(x)}}
$$
is strictly convex on $(0,\infty)$.
\end{lem}

\begin{proof}
Define $\Lambda(x)=\ln\frac{f(x)}{f(x+a)}$ for $x\in(0,\infty)$.
Since $\frac{f'(x)}{f(x)}=-\frac1{H(x)}$, we have
$$
\Lambda(x)=\ln f(x)-\ln f(x+a)
=\int_{x+a}^{x}\frac{f'(t)}{f(t)}\td t
=
\int_x^{x+a}\frac{\td t}{H(t)}>0.
$$
Consequently, it follows that
$$
\Lambda'(x)=\frac1{H(x+a)}-\frac1{H(x)}<0
$$
and
$$
\Lambda''(x)
=\frac{H'(x)}{H^2(x)}-\frac{H'(x+a)}{H^2(x+a)}.
$$
Therefore, we acquire
\begin{equation}\label{eq:lambda-bound}
\begin{aligned}
\frac{\Lambda''(x)}{[\Lambda'(x)]^2}
&=\frac{H'(x)H^2(x+a)-H'(x+a)H^2(x)}{[H(x+a)-H(x)]^2}\\
&\leq \frac{H'(x)[H^2(x+a)-H^2(x)]}{[H(x+a)-H(x)]^2}\\
&=H'(x)\frac{H(x+a)+H(x)}{H(x+a)-H(x)}\\
&=H'(x)\coth\frac{L(x)}{2}, \quad x\in(0,\infty),
\end{aligned}
\end{equation}
where
$$
L(x)=\ln\frac{H(x+a)}{H(x)}>0, \quad x\in(0,\infty).
$$
Moreover, it is easy to see that
$$
L(x)=\int_x^{x+a}\frac{H'(t)}{H(t)}\td t
\geq H'(x)\int_x^{x+a}\frac{\td t}{H(t)}
=H'(x)\Lambda(x).
$$
Hence, we arrive at
\begin{equation}\label{LL-ineq}
\frac{\Lambda(x)}{2}\leq\frac{L(x)}{2H'(x)}, \quad x\in(0,\infty).
\end{equation}
\par
Let
$$
G(z)=z\coth z>0, \quad z\in(0,\infty).
$$
Direct differentiation gives
$$
G'(z)=\frac{\sinh z\cosh z-z}{\sinh^2z}
$$
and
$$
(\sinh z\cosh z-z)'=2\sinh^2z>0.
$$
This implies that $G'(z)>0$ and then $G(z)$ is strictly increasing for $z\in(0,\infty)$.

Since $0<H'(x)<1$, it follows that
$$
\frac{L(x)}{2H'(x)}>\frac{L(x)}{2}.
$$
The strict increase of $G(z)$ results in
$$
\coth\frac{L(x)}{2H'(x)}>H'(x)\coth\frac{L(x)}{2}.
$$
Since $\coth z$ is strictly decreasing on $(0,\infty)$ and the inequality~\eqref{LL-ineq}, we obtain
$$
\coth\frac{\Lambda(x)}{2}
\geq
\coth\frac{L(x)}{2H'(x)}
>
H'(x)\coth\frac{L(x)}{2}.
$$
Combining this with~\eqref{eq:lambda-bound} leads to the inequality
\begin{equation}\label{eq:key-inequality}
\frac{\Lambda''(x)}{[\Lambda'(x)]^2}<\coth\frac{\Lambda(x)}{2}, \quad x\in(0,\infty).
\end{equation}
\par
Now assume that
$$
q(x)=\te^{-\Lambda(x)}=\frac{f(x+a)}{f(x)}, \quad x\in(0,\infty).
$$
Then $0<q(x)<1$ and
$$
F_a(x)=\frac1{1-q(x)}.
$$
Furthermore, we have
$$
q'(x)=-q(x)\Lambda'(x)\quad\text{and}\quad
q''(x)=q(x)\bigl([\Lambda'(x)]^2-\Lambda''(x)\bigr).
$$
Thus, the second derivative
$$
F_a''(x)=\frac{q''(x)}{[1-q(x)]^2}+\frac{2[q'(x)]^2}{[1-q(x)]^3}
=
\frac{q(x)}{[1-q(x)]^2}
\left[
\frac{1+q(x)}{1-q(x)}[\Lambda'(x)]^2-\Lambda''(x)
\right].
$$
Since
$$
\frac{1+q(x)}{1-q(x)}
=
\coth\frac{\Lambda(x)}{2},
$$
we obtain
$$
F_a''(x)
=
\frac{\te^{-\Lambda(x)}[\Lambda'(x)]^2}{[1-\te^{-\Lambda(x)}]^2}
\biggl[\coth\frac{\Lambda(x)}{2}-\frac{\Lambda''(x)}{[\Lambda'(x)]^2}\biggr].
$$
The bracket is strictly positive by~\eqref{eq:key-inequality}.
Therefore, it follows that $F_a''(x)>0$. The proof of Lemma~\ref{lem:general} is thus complete.
\end{proof}

\end{document}
